\documentclass[conference]{IEEEtran}
\usepackage{cite}
\usepackage{amsmath,amssymb,amsfonts}
\usepackage{graphicx}
\usepackage{textcomp}
\usepackage{xcolor}
\begin{document}

\title{Language-Guided Autonomous Navigation for a Simulated Husky Using Qwen Vision-Language Models}

\author{\IEEEauthorblockN{1\textsuperscript{st} Lev Saunders}
\IEEEauthorblockA{\textit{Computer Science} \\
\textit{Drexel University}\\
Philadelphia, PA \\
las564@drexel.edu}}

\maketitle

\begin{abstract}
This project will extend the existing ROS autonomous-navigation simulation of a Clearpath Husky for HW1 by adding language and vision-guided decision making. An off-the-shelf Qwen vision-language model (VLM) will receive a natural-language navigation instruction together with images from the Husky's simulated camera. The model will select navigation actions for the robot to perform. A ROS 2 interface node will translate that decision into a short-range navigation goal for Nav2, which will remain responsible for path planning and low-level motion. The final implementation will demonstrate closed-loop language-guided navigation in a simulated environment using tasks such as approaching a visible object or passing through a doorway.
\end{abstract}


\section{Project Motivation and Objective}
In the HW1, I used the Clearpath Husky sample configuration provided for simulation~\cite{clearpathyaml} in order to implement autonomous navigation using ROS 2 Jazzy and Nav2. The robot can navigate to geometric goals while using its range sensors for localization and obstacle avoidance. However, conventional Nav2 goals are normally specified as map coordinates and do not by themselves allow the robot to interpret commands such as ``go through the doorway'' or ``approach the blue chair.'' The objective of this project is to add a lightweight semantic decision-making layer that connects natural-language instructions and visual observations to the existing navigation system.

An available Qwen VLM will be used directly for visual and language reasoning. Qwen3-VL provides multimodal reasoning, spatial perception, and visual grounding capabilities that are useful for connecting visible scene content to simple high-level navigation decisions~\cite{qwen3vl}. The project will use these capabilities without training or fine-tuning the model, keeping the implementation focused on integration with the existing ROS 2 navigation stack.

\section{Proposed System}
The existing Husky simulation will remain largely unchanged. The robot configuration includes an Intel RealSense D435 RGB-D camera as well as 2D and 3D lidar sensors. Nav2 will continue to perform geometric navigation and collision-aware motion planning. The new component will be a ROS 2 node that interfaces between the simulated camera, a user instruction, the Qwen model, and Nav2.

At each decision step, the node will capture the current RGB camera image and combine it with a persistent user instruction. For example, the user may provide the command ``go through the doorway.'' The image and instruction will be passed to Qwen using a constrained prompt. Rather than asking the general-purpose VLM to generate precise metric trajectories, the model will choose from a small action vocabulary such as \texttt{FORWARD}, \texttt{FORWARD\_LEFT}, \texttt{FORWARD\_RIGHT}, \texttt{LEFT}, \texttt{RIGHT}, and \texttt{STOP}. The output will be constrained to a machine-readable format so that invalid responses can be rejected.

The selected action will then be converted into a short local goal relative to the Husky. For example, \texttt{FORWARD} may correspond to a goal approximately one meter ahead, while \texttt{FORWARD\_LEFT} will produce a nearby goal ahead and to the left. The goal will then be transformed into the appropriate coordinate frame and sent to the Nav2 \texttt{NavigateToPose} action server. Nav2, rather than the VLM, will determine the safe path and generate velocity commands.

After the Husky reaches the local goal, the system will acquire a new camera observation and query Qwen again. This creates a simple closed-loop process of observation, semantic decision, local navigation, and re-observation. The loop terminates when Qwen selects \texttt{STOP}, the instruction is completed, or a maximum number of steps is reached. This separation keeps the learned component focused on high-level decisions while relying on the existing navigation stack for reliable robot control.

\section{Implementation Plan}
Development will proceed incrementally. First, Qwen inference will be tested locally on static screenshots from the Gazebo simulation to determine a prompt and output format that reliably produces one of the permitted actions. A smaller Qwen3-VL checkpoint will be preferred so that inference can run on the available RTX 3070; image resolution and model quantization may be reduced if necessary.

Next, a ROS 2 perception/decision node will subscribe to the camera topic and accept a text instruction. The node will package the latest image and instruction for model inference and parse the returned action. A separate navigation interface will translate each discrete action into a short relative pose and submit it to Nav2. Safety checks will prevent malformed model responses or unreasonable goals from being executed.

Finally, these components will be combined into the closed-loop navigation agent. The initial target is a small set of visually grounded tasks in one simulated indoor environment. Examples of potential demonstrations are ``go through the doorway,'' ``approach the blue chair,'' and ``navigate toward the desk.''

\section{Evaluation and Deliverables}
The goal of the project is to demonstrate a working implementation. For each supported instruction, the Husky should interpret the camera scene, produce valid decisions, send goals to Nav2, and make progress toward the requested target without bypassing Nav2's obstacle avoidance. Testing will record whether each task is completed, the number of VLM decisions required, and approximate inference latency. Failure cases, such as incorrect object identification, oscillating decisions, or premature \texttt{STOP} actions, will also be documented.

The final deliverables will include the ROS 2 nodes used for Qwen inference and Nav2 integration, launch/configuration files required to reproduce the simulated demonstration, and a short video or live demonstration of the Husky responding to natural-language navigation commands. The final report will describe the software architecture, implementation decisions, limitations, and observed behavior. The intended contribution is a manageable demonstration of how a general-purpose vision-language model can serve as a semantic decision layer on top of a conventional autonomous navigation stack.

\begin{thebibliography}{00}
\bibitem{clearpathyaml} Clearpath Robotics, ``Robot YAML,'' Clearpath Robotics Documentation, ROS 2 Jazzy. [Online]. Available: docs.clearpathrobotics.com/docs/ros/config/yaml/overview/. Accessed: Oct. 25, 2026.
\bibitem{qwen3vl} S. Bai et al., ``Qwen3-VL Technical Report,'' arXiv:2511.21631, 2025.
\end{thebibliography}

\end{document}